\clearpage
\onecolumn
\thispagestyle{empty}
\vspace*{\fill}
\begin{center}
	{\color{radblue}\rule{\linewidth}{2pt}}\par
	\vspace{6mm}
	{\Huge\bfseries\color{radblue} Medición de huesos largos}\par
	\vspace{6mm}
	{\color{radblue}\rule{\linewidth}{2pt}}
\end{center}
\vspace*{\fill}
\clearpage
\twocolumn
\pagestyle{contenido}


%%%%%%%%%%%%%%%%%%%%%%%%%%%%%%%%%%%%%%%%%%%%%%%%%%%%%%%%%%%
%  SECCIÓN: Pelvis
%  Archivo: secciones/pelvis.tex
%  Sin preámbulo ni \begin{document} — lo maneja main.tex
	%%%%%%%%%%%%%%%%%%%%%%%%%%%%%%%%%%%%%%%%%%%%%%%%%%%%%%%%%%%
	
	\section{Medición de huesos largos}
	
	\subsection{Pelvimetría - AP de Pie}
	
	\subsubsection*{1. Identificación}
	
	\textbf{Tipo:} Proyección funcional. Estudio métrico de la pelvis orientado a la evaluación de ángulos y medidas estructurales del anillo pélvico.
	
	\textbf{Indicaciones clínicas principales:}
	\begin{itemize}
		\item Evaluación pré-quirúrgica o pré-obstétrica del diámetro pélvico
		\item Medición de dismetría de caderas y valoración de ángulos femorales
		\item Seguimiento de displasia del desarrollo de cadera
	\end{itemize}
	
	\subsubsection*{2. Técnica Base}
	
	\textbf{Receptor de imagen:} 35$\times$43 cm. Orientación \textbf{transversal (apaisado)}. El receptor debe ser cuadriculado especial para permitir las mediciones métricas. En ausencia del chasis cuadriculado, puede emplearse software de medición en el negatoscopio o sistema PACS.
	
	\textbf{DFRI:} 100 cm. Distancia estándar para técnica convencional de pelvis, permite reproducibilidad de las mediciones y minimiza la magnificación geométrica.
	
	\textbf{Técnica:} Intermedia, equivalente a la técnica estándar de pelvis. Permite visualizar partes blandas y estructuras óseas en equilibrio adecuado de contraste para realizar las mediciones.
	
	\textbf{Bucky:} Bucky mural con grilla. Se utiliza bucky mural ya que el paciente se posiciona de pie.
	
	\subsubsection*{3. Posicionamiento}
	
	\paragraph{Paciente}
	Paciente de pie, descalzo, de frente al receptor. La posición descalza es fundamental para evitar la influencia del calzado en la nivelación de la pelvis y garantizar la reproducibilidad de las mediciones.
	
	% TODO: \usepackage{graphicx} required
	\begin{figure}[h]
		\centering
		\includegraphics[width=0.7\linewidth]{imagenes/medi1}
		\caption{}
		\label{fig:medi1}
	\end{figure}
	
	
	\paragraph{Región de interés}
	Miembros inferiores con rotación interna para desplegar los cuellos femorales en toda su longitud, permitiendo su correcta valoración métrica. Esta rotación contrarrestará la anteversión femoral fisiológica.
	
	\paragraph{Rayo Central}
	Perpendicular al plano del receptor. Punto de entrada: \textbf{punto medio entre ambas espinas ilíacas anterosuperiores (EIAS)}.
	
	\paragraph{Receptor de imagen}
	Centrado al punto medio entre las EIAS. La región a incluir abarca desde las crestas ilíacas superiormente hasta la parte proximal de ambos fémures inferiormente, incluyendo la totalidad del anillo pélvico.
	
	\paragraph{Respiración}
	Respiración normal, sin apnea requerida. La región pélvica no se ve afectada por los movimientos respiratorios en un estudio en bipedestación.
	
	\subsubsection*{4. Anatomía Visible}
	
	Anillo pélvico completo con ambas alas ilíacas, sacro y cóccix. Articulaciones sacroilíacas simétricas. Cabezas femorales, cuellos femorales desplegados por la rotación interna, trocánteres mayores y menores. Articulaciones coxofemorales bilaterales. Sínfisis del pubis al centro.
	
	Las mediciones se realizan sobre el receptor cuadriculado o mediante software, trazando las siguientes referencias:
	\begin{itemize}
		\item \textbf{Línea A:} paralela a la línea de eburnación
		\item \textbf{Líneas B:} paralelas a la plomada a través de las cabezas femorales
		\item \textbf{Línea C:} perpendicular a la línea de la plomada
		\item \textbf{Valor D:} dismetría medida entre ambas cabezas femorales (Tilley, P.; 1966. Irvin, R., 1991). La dismetría puede medirse también con caliper digital en negatoscopio.
	\end{itemize}
	
	% TODO: \usepackage{graphicx} required
	\begin{figure}
		\centering
		\includegraphics[width=0.7\linewidth]{imagenes/medi3}
		\caption{}
		\label{fig:medi3}
	\end{figure}
	
	
	\subsubsection*{5. Criterios de Evaluación}
	
	\textbf{Posicionamiento correcto:} Ambas espinas ilíacas anterosuperiores equidistantes del borde lateral del receptor. Cuellos femorales visibles en su totalidad sin acortamiento. Sacro proyectado en la línea media.
	
	\textbf{Ausencia de rotación:} Simetría de las alas ilíacas bilaterales. Agujeros obturadores de igual forma y tamaño a ambos lados. Articulaciones sacroilíacas simétricas.
	
	\textbf{Penetración y contraste:} Visualización simultánea de partes blandas pélvicas y estructuras óseas. Cuerpos vertebrales lumbares bajos visibles con densidad adecuada. Trabéculas del hueso esponjoso de cabezas femorales perceptibles para permitir las mediciones métricas.
	
	\subsubsection*{6. Puntos Críticos}
	
	\textbf{Calzado no retirado:} Introduce una diferencia artificial en la altura pélvica que invalida la medición de dismetría. Siempre verificar que el paciente esté descalzo antes de la exposición.
	
	\textbf{Rotación interna insuficiente de MMII:} Los cuellos femorales se proyectan acortados o superpuestos con los trocánteres mayores, imposibilitando la correcta medición del eje femoral. Corregir orientando los pies con las puntas hacia adentro.
	
	\textbf{Adaptaciones:} En pacientes con limitación de la rotación interna (prótesis, artrosis severa), documentar la posición real utilizada ya que afectará la interpretación de las mediciones. En pacientes que no pueden mantenerse de pie, consultar con el médico la viabilidad de técnica alternativa.
	
	\textbf{Trucos:} Si no se dispone del receptor cuadriculado especial, el estudio puede realizarse con receptor estándar y las mediciones se efectúan posteriormente mediante software en la estación de trabajo o PACS. La posición de pie es inherente a la naturaleza funcional del estudio, ya que reproduce la carga axial real sobre la pelvis.
	
	\onecolumn
	\clearpage
	
	% TODO: \usepackage{graphicx} required
	\begin{figure}
		\centering
		\includegraphics[width=0.7\linewidth]{imagenes/medi2}
		\caption{}
		\label{fig:medi2}
	\end{figure}
	
	\twocolumn
	\clearpage
	%----------------------------------------------------------
	%  Para agregar otro enfoque duplicá el bloque
	%  desde \subsection hasta \textbf{Trucos}
	%----------------------------------------------------------
	
%%%%%%%%%%%%%%%%%%%%%%%%%%%%%%%%%%%%%%%%%%%%%%%%%%%%%%%%%%%
%  SECCIÓN: Miembro Inferior — Medición
%  Archivo: secciones/green_ortorradiografia.tex
%  Sin preámbulo ni \begin{document} — lo maneja main.tex
	%%%%%%%%%%%%%%%%%%%%%%%%%%%%%%%%%%%%%%%%%%%%%%%%%%%%%%%%%%%
	
	\section{Medición de huesos largos}
	
	\subsection{Método de Green - Ortorradiografía (AP Bilateral Comparativa)}
	
	\subsubsection*{1. Identificación}
	
	\textbf{Tipo:} Estudio bilateral comparativo. Comprende tres enfoques secuenciales sobre el mismo receptor: articulación de la cadera, rodilla y tobillo (tibiotarsiana).
	
	\textbf{Indicaciones clínicas principales:}
	\begin{itemize}
		\item Dismetría funcional o estructural de miembros inferiores
		\item Escoliosis o desviación de columna por diferencia de longitud de miembros
		\item Control de crecimiento (hipometría o hipermetría)
	\end{itemize}
	
	\subsubsection*{2. Técnica Base}
	
	\textbf{Receptor de imagen:} 35$\times$43 cm. Orientación \textbf{apaisada (transversal)}, dividido en tres franjas de 14 cm aproximadamente, una por articulación.
	
	\textbf{DFRI:} 100 cm. Distancia estándar para este estudio; su cumplimiento estricto es condición indispensable para que las mediciones sean fidedignas.
	
	\textbf{Técnica:} Se calculan \textbf{tres técnicas independientes}, una por articulación.
	\begin{itemize}
		\item \textbf{Cadera:} técnica intermedia, ya que el RC ingresa al PMS y no está focalizado en la cadera; al ser un estudio bilateral se atraviesa más tejido y se genera mayor radiación dispersa. Esta técnica actúa como base, aportando densidad uniforme para ambas articulaciones.
		\item \textbf{Rodilla:} técnica homogeneizada respecto al disparo de cadera.
		\item \textbf{Tobillo:} técnica homogeneizada respecto al disparo de cadera.
	\end{itemize}
	El disparo de cadera permite visualizar desde la densidad alta del fémur proximal hasta las partes blandas y el trabeculado óseo a nivel distal.
	
	\textbf{Bucky:} Bucky de mesa con grilla. Se desplaza tanto el tubo de rayos como el receptor de imagen entre cada disparo.
	
	\subsubsection*{3. Posicionamiento}
	
	\paragraph{Paciente}
	Decúbito supino (proyección AP), inmovilizado durante las tres tomas. Las piernas se cinchan para mantener la posición. Se solicita al paciente que se relaje y permanezca flojo. Se le debe retirar el pantalón o subirlo lo máximo posible. Si presenta prótesis, se tracciona desde los hombros; de lo contrario, desde los pies. La comunicación constante con el paciente durante el estudio es fundamental para garantizar su colaboración y la fidelidad de las mediciones.
	
	\begin{figure}[h]
		\centering
		\includegraphics[width=0.7\linewidth]{imagenes/medi3}
		\caption{}
		\label{fig:medi3}
	\end{figure}
	
	
	\paragraph{Región de interés}
	Pelvis sin rotación, con las EIAS equidistantes al receptor de imagen. Miembros inferiores extendidos con rotación medial de 10--15°.
	
	\paragraph{Rayo Central}
	\begin{itemize}
		\item \textbf{Cadera:} RC perpendicular al receptor, punto de entrada 2,5 cm proximal a la sínfisis del pubis.
		\item \textbf{Rodilla:} RC perpendicular al receptor, a la altura del vértice de la rótula.
		\item \textbf{Tobillo:} RC perpendicular al receptor, a la altura de los maléolos.
	\end{itemize}
	
	\paragraph{Receptor de imagen}
	Centrado de modo que cada articulación quede incluida en su franja correspondiente del receptor. Las tres articulaciones se registran en el mismo chasis. La región a incluir comprende articulaciones de cadera, rodillas y tobillos (tibiotarsianas).
	
	\paragraph{Respiración}
	Respiración normal. No se requiere apnea. Lo determinante para la validez del estudio es la inmovilidad del paciente durante cada disparo.
	
	\subsubsection*{4. Anatomía Visible}
	
	Articulaciones de cadera bilateral, rodillas y tobillos en una sola imagen. Fémur proximal con alta densidad ósea, trabeculado y partes blandas a nivel distal. En pacientes pediátricos se visualizan las líneas de crecimiento, lo cual es un hallazgo de relevancia clínica. El estudio no presenta magnificación: es fidedigno a la estructura anatómica real, a diferencia de la telerradiografía.
	
	\subsubsection*{5. Criterios de Evaluación}
	
	\textbf{Posicionamiento correcto:} Cada articulación centrada en su franja del receptor; pelvis sin rotación con EIAS simétricas; miembros inferiores con rotación medial uniforme de 10--15°.
	
	\textbf{Ausencia de rotación:} Simetría bilateral de las articulaciones de cadera, rodillas y tobillos; fémures sin rotación visible.
	
	\textbf{Penetración y contraste:} Densidad adecuada en cadera (fémur proximal y partes blandas visibles simultáneamente); rodilla y tobillo con escala de grises homogeneizada respecto al disparo de cadera. Las tres articulaciones deben mostrar detalle trabecular.
	
	\subsubsection*{6. Puntos Críticos}
	
	\textbf{Movimiento del paciente entre disparos:} Invalida las mediciones de dismetría; el estudio deja de ser fidedigno. Corrección: inmovilizar con cinchas, explicar el procedimiento, hablarle al paciente antes de cada disparo indicándole que mantenga la posición y mirarlo constantemente durante el estudio.
	
	\textbf{Falta de protección gonadal:} Debe colocarse el protector gonadal sin que interfiera en el área de interés, especialmente en niños, donde también se deben visualizar las líneas de crecimiento.
	
	\textbf{Adaptaciones:} En pacientes con prótesis, la tracción se realiza desde los hombros en lugar de los pies. En pacientes pediátricos, verificar visibilidad de líneas de crecimiento en cada disparo.
	
	\textbf{Trucos:} Este estudio no muestra el hueso en su longitud completa, lo que impide visualizar áreas deformadas o patológicas de las diáfisis femorotibiales (limitación compensada por la telerradiografía). La diferencia de longitud entre miembros debe ser mínima para que el estudio sea concluyente. Recordar que es la colimación estrecha y el desplazamiento coordinado de tubo y receptor lo que permite los tres enfoques en un único chasis.
	
	\onecolumn
	\clearpage
	
	\begin{figure}
		\centering
		\includegraphics[width=0.7\linewidth]{imagenes/medi4}
		\caption{Una pierna}
		\label{fig:medi4}
	\end{figure}
	
	\begin{figure}
		\centering
		\includegraphics[width=0.7\linewidth]{imagenes/medi5}
		\caption{Ambas piernas}
		\label{fig:medi5}
	\end{figure}
	
	
	\twocolumn
	\clearpage
	%----------------------------------------------------------
	%  Para agregar otro enfoque duplicá el bloque
	%  desde \subsection hasta \textbf{Trucos}
	%----------------------------------------------------------
	
	
%%%%%%%%%%%%%%%%%%%%%%%%%%%%%%%%%%%%%%%%%%%%%%%%%%%%%%%%%%%
%  SECCIÓN: Goniometría
%  Archivo: secciones/goniometria.tex
%  Sin preámbulo ni \begin{document} — lo maneja main.tex
	%%%%%%%%%%%%%%%%%%%%%%%%%%%%%%%%%%%%%%%%%%%%%%%%%%%%%%%%%%%
	
	\section{Medición de huesos largos}
	
	\subsection{Goniometría - Telerradiografía de Miembros Inferiores}
	
	\subsubsection*{1. Identificación}
	
	\textbf{Tipo:} Estudio funcional comparativo bilateral o monopodalico. Puede realizarse de forma bilateral o monopodalica según indicación clínica.
	
	\textbf{Indicaciones clínicas principales:}
	\begin{itemize}
		\item Evaluación y medición de deformidades angulares del miembro inferior (genu varo, genu valgo)
		\item Planificación prequirúrgica en ortopedia y traumatología
		\item Seguimiento evolutivo de correcciones terapéuticas y patologías degenerativas
		\item Evaluación de secuelas post-traumáticas en medicina laboral y legal (determinación de incapacidad)
	\end{itemize}
	
	\subsubsection*{2. Técnica Base}
	
	\textbf{Receptor de imagen:} 35 $\times$ 43 cm. Orientación \textbf{vertical} en posición mural; en equipos convencionales el receptor se divide en tres sectores longitudinales de 14 cm cada uno para la técnica de tres exposiciones. En niños puede utilizarse un solo receptor según el tamaño del paciente.
	
	\textbf{DFRI:} 180 cm mínimo; óptimo a 200--250 cm. La mayor distancia reduce la magnificación inherente al estudio, aunque persiste un error aproximado del 15\% que puede paliarse con el uso de reglas radio-opacas de longitud conocida colocadas a la misma distancia que los huesos a medir.
	
	\textbf{Técnica:} Intermedia y homogenizada, adaptada a cada segmento anatómico según su grosor. Los valores de kVp se ajustan por región: cadera 75 kVp, rodilla 62 kVp, tobillo 54 kVp. Los mAs son variables y se calculan en función del grosor de la región a exponer. Al aumentar la DFRI a 2 metros se aplica compensación por ley del cuadrado inverso (factor $\times$4); técnica intermedia de referencia: 140 mAs a 2 metros.
	
	\textbf{Bucky:} Bucky mural con grilla. El paciente se posiciona con la espalda apoyada contra el bucky mural durante todo el estudio.
	
	\subsubsection*{3. Posicionamiento}
	
	\paragraph{Paciente}
	De pie, descalzo, sobre un cajón de madera elevador, con la espalda apoyada contra el bucky mural. El peso corporal debe distribuirse de forma simétrica y equitativa sobre ambos miembros. La tracción para alinear el miembro se realiza desde los miembros inferiores; en pacientes con prótesis de cadera se utiliza tracción axilar para evitar complicaciones.
	
	\begin{figure}[h]
		\centering
		\includegraphics[width=0.7\linewidth]{imagenes/medi6}
		\caption{En esta imagen se usa un chasis propio para telerradiografia}
		\label{fig:medi6}
	\end{figure}
	
	
	\paragraph{Región de interés}
	Leve rotación interna de los miembros inferiores para desplegar el cuello femoral y estandarizar la posición. Las rótulas deben quedar orientadas al frente como referencia de ausencia de rotación. Ambos miembros deben mantenerse en posición comparable y simétrica durante todo el estudio.
	
	\paragraph{Rayo Central}
	Perpendicular al plano del receptor. El RC se centra a la altura de las rodillas como referencia general del estudio. En la técnica de tres exposiciones el centrado varía por segmento: primera exposición en la articulación de la cadera (2 cm por encima de la sínfisis púbica), segunda en el vértice de la rótula (articulación de la rodilla), tercera en la articulación tibiotarsiana (tobillo).
	
	\paragraph{Receptor de imagen}
	El receptor se centra de forma secuencial en cada una de las tres articulaciones principales. Entre exposiciones se desplaza únicamente el receptor (bucky), mientras el tubo y el paciente permanecen absolutamente inmóviles. Las imágenes deben solaparse aproximadamente 2 cm para garantizar la continuidad anatómica y permitir la unión digital posterior. La colimación debe mantenerse totalmente abierta (apertura máxima) para cubrir desde la articulación coxofemoral hasta los maléolos.
	

	
	\paragraph{Respiración}
	Respiración normal durante todo el estudio. No se requiere apnea. La inmovilización estricta del paciente es el factor crítico, no la fase respiratoria.
	
	\subsubsection*{4. Anatomía Visible}
	
	El estudio debe incluir las tres articulaciones principales del miembro inferior en su totalidad: articulación coxofemoral (cadera), articulación fémoro-tibial (rodilla) y articulación tibiotarsiana (tobillo). Se visualizan fémur completo, tibia y peroné en toda su extensión. La orientación de cada articulación se evalúa mediante líneas de referencia específicas: en cadera, línea trazada de la punta del trocánter mayor al centro de la cabeza femoral; en rodilla, dos trazos independientes sobre el extremo distal del fémur y el extremo proximal de la tibia, con un ángulo convergente normal de 3° entre ambos; en tobillo, línea trazada a nivel del extremo distal de la tibia. A partir de estas referencias se traza el eje mecánico del miembro inferior, línea imaginaria que une el centro de la cabeza femoral con el centro de la articulación tibiotarsiana y que en condiciones normales pasa por el centro de la rodilla con una desviación máxima de 8°.
	
	\subsubsection*{5. Criterios de Evaluación}
	
	\textbf{Posicionamiento correcto:} Las tres articulaciones principales deben estar completamente incluidas y visibles. Las imágenes parciales deben solaparse correctamente para permitir la reconstrucción del miembro completo. La imagen final debe comprender desde la articulación coxofemoral hasta los maléolos sin discontinuidad anatómica.
	
	\textbf{Ausencia de rotación:} Las rótulas deben proyectarse al frente, centradas entre los cóndilos femorales. La simetría bilateral en estudios comparativos confirma el posicionamiento estándar. El cuello femoral debe verse desplegado gracias a la rotación interna leve.
	
	\textbf{Penetración y contraste:} Cada segmento debe presentar densidad radiográfica adecuada a su grosor, lograda mediante la adaptación individual de los mAs. La técnica homogenizada permite visualizar correctamente las estructuras óseas en los tres niveles sin sobreexposición ni subexposición en ninguno de los sectores.
	
	\subsubsection*{6. Puntos Críticos}
	
	\textbf{Movimiento del paciente entre exposiciones:} Invalida el estudio completo ya que impide la alineación de las imágenes y la medición del eje mecánico. Se debe comunicar con claridad al paciente la necesidad de inmovilización estricta durante toda la adquisición, explicitando que el estudio es prolongado.
	
	\textbf{Asimetría de carga o posición:} Una distribución desigual del peso o una posición no estandarizada altera los ángulos medidos y produce resultados no reproducibles. Verificar carga simétrica y rotación interna leve antes de iniciar las exposiciones.
	
	\textbf{Adaptaciones:} En pacientes con prótesis de cadera sustituir la tracción desde miembros inferiores por tracción axilar. En pacientes pediátricos puede utilizarse un único receptor si el tamaño del paciente lo permite, y se deben adaptar los parámetros técnicos. Ante limitaciones de colaboración considerar si el estudio completo es clínicamente necesario.
	
	\textbf{Trucos:} Las medidas obtenidas son comparativas, no absolutas; el traumatólogo realiza las mediciones sobre la imagen final reconstruida. La magnificación inherente (aprox. 15\%) no se elimina pero puede paliarse con reglas radio-opacas de longitud conocida colocadas a la misma distancia que los huesos. En equipos con espinómetro digital el receptor se desplaza durante la exposición permitiendo captura continua del miembro en una sola imagen, lo que elimina el problema del solapamiento manual.
	
		\begin{figure}[h]
		\centering
		\includegraphics[width=0.7\linewidth]{imagenes/medi7}
		\caption{División del chasis}
		\label{fig:medi7}
	\end{figure}
	
	\onecolumn
	\clearpage
	
	\begin{figure}
		\centering
		\includegraphics[width=0.5
		\linewidth]{imagenes/medi8}
		\caption{}
		\label{fig:medi8}
	\end{figure}
	
	
	\twocolumn
	\clearpage
	%----------------------------------------------------------
	%  Para agregar otro enfoque duplicá el bloque
	%  desde \subsection hasta \textbf{Trucos}
	%----------------------------------------------------------
	
	
